\documentclass[12pt]{article}
\usepackage[utf8]{inputenc}
\usepackage[T1]{fontenc}
\usepackage{hyperref}
\title{Waama Language Resources for K-12 Curriculum Development}
\author{Compiled by Claude Sonnet 5.5 (claude-sonnet-5-5) for Andr\'as Kornai}
\date{}
\begin{document}
\maketitle

\section{Introduction}
Waama (also called Yoabu; ISO 639-3: wwa) is a Gur (Niger-Congo) language of northern Benin, in the Atacora department. Speaker estimates vary widely: about 50,000 (Wikipedia, from 2000), about 127,000 (a linguistic data entry) and 141,000 (ASJP). It is a recognised minority language in Benin. It has two tones, and the orthography marks tone with a grave accent on certain letters. In searches I found a dictionary and one grammar, but no school textbooks.

\section{Primary Educational Resources}
No primers or textbooks were found in searches. The Waama dictionary states that it aims to give the language a standard written form and to serve teachers and literacy workers. That suggests literacy work exists, but I saw no URL for any literacy material.

\section{Linguistic and Reference Materials}
\begin{itemize}
\item Waama--French dictionary, produced by SIL Benin and the Centre National de Linguistique Appliqu\'ee (CENALA) and following CENALA orthographic standards. It is hosted on Webonary: \url{https://www.webonary.work/waama/overview/forward/?lang=en}. Entry conventions are explained at \url{https://www.webonary.work/waama/?p=146827}. Verbs are given in three forms (imperative, perfective, progressive), and nouns in singular and plural.
\item Ursula Bremicker (1990), a French-language dissertation describing Waama phonology and grammar (ANRT, Lille). Bibliographic record: \url{https://glottolog.org/resource/reference/id/32302}.
\item General overview: \url{https://en.wikipedia.org/wiki/Waama_language}.
\end{itemize}

\section{Religious and Community Resources}
No scripture or community-organisation resources for Waama appeared in the search results.

\section{Digital and Audio Resources}
The online Webonary dictionary is the only digital resource found. No audio or video resources appeared in the results.

\section{Curriculum Development Context}
Waama is spoken by a sizeable community, and its orthography has been standardised by CENALA. This gives a good basis for literacy materials. The dictionary and the Bremicker grammar both support teacher training. The missing pieces are graded readers, primers and mother-tongue subject materials. Benin's official language of instruction is French, so any curriculum would probably be bilingual French--Waama.

\end{document}